% TOP_PROBLEM: {"id":30007587,"problem_number":"LOCAL-30007587","title":"Motivated beliefs in economic decisions","category_id":21,"category":{"id":21,"name":"economics","display_name":"Economics","description":"Open economic theory statements, published research agendas, and proposed scoped specifications; question provenance and historical priority are qualified per record.","slug":"economics","order_index":21,"created_at":"2026-10-08T00:00:00Z"},"status":"provenance_pending","external_url":null,"legacy_ids":[],"published":false,"statement":"Identify how incentives for desirable conclusions change economic beliefs and actions under the specified randomized design and measurement assumptions.","economics":{"original_id":76,"field":"Behavioral and experimental economics","kind":"identification","formalization_basis":"proposed_specification","source_status":"not_verified","priority_status":"not_established","priority_note":"The source supplies evidence in stock-return experiments, not verified priority or an explicit statement of the broad saving and self-image question.","earliest_verified_reference":null,"current_reference":{"authors":["Carlos Cueva","Iñigo Iturbe-Ormaetxe"],"title":"Motivated beliefs about stock returns","year":2025,"url":"https://doi.org/10.1016/j.jbankfin.2025.107510","doi":"10.1016/j.jbankfin.2025.107510","locator":"Abstract, corroborated by author research page https://sites.google.com/site/carloscuevaherrero/home/research. Full paper was not accessible in this audit.","evidence_summary":"The authors report experimental effects of stock purchase on optimistic expectations after losses and a relationship with reluctance to sell. The source does not verify the broader saving, self-image, persistence, and prevalence question as a stated agenda."},"formalization_note":"Proposed causal research design extending a supporting experiment. Full-paper agenda and persistence/prevalence claims could not be verified."},"statusQualification":"Provenance pending: an explicit published statement of this exact open question was not verified. This is a catalogue-authored candidate specification, not a certified open conjecture. Proposed causal research design extending a supporting experiment. Full-paper agenda and persistence/prevalence claims could not be verified.","statusReviewedAt":"2026-10-08"}
% FIRST_POSTED: 2026-10-08
% ENTRY_KIND: statement_only
% !TeX program = lualatex
\documentclass[11pt]{article}
\usepackage{fontspec}
\usepackage[a4paper,margin=25mm]{geometry}
\usepackage{amsmath,amssymb,mathtools}
\usepackage{xurl}
\usepackage[hidelinks]{hyperref}
\newcommand{\sref}[2]{\hyperlink{source:#1:#2}{[S#2]}}
\setlength{\emergencystretch}{3em}
\begin{document}
% problemId:30007587
\section*{Motivated beliefs in economic decisions}\label{30007587}
\subsection{Definitions and mathematical statement}
Fix an economic decision with uncertain monetary return \(R\) and objectively specified public information \(I\). Let \(Z_i\in\{0,1\}\) randomly assign an ownership or self-image treatment intended to change the desirability of a favorable return without changing available information. At horizon \(t\), observe an incentivized subjective return forecast \(B_i(Z_i,t)\) and an economic action \(A_i(Z_i,t)\), such as a portfolio share. The belief and action effects in target population \(P\) are
\[\tau_B(t)=\mathbb E_P[B_i(1,t)-B_i(0,t)],\qquad \tau_A(t)=\mathbb E_P[A_i(1,t)-A_i(0,t)].\]
Specify the information benchmark \(b^*(I)=\mathbb E[R\mid I]\) and examine whether the treatment changes forecast errors relative to it. A motivation interpretation requires ruling out treatment-induced information, altered forecast incentives, or rational inferences about the experiment. To determine whether distorted beliefs cause the action effect, additionally randomize a credible belief-correction intervention or justify an explicit mediation model; correlation between \(B_i\) and \(A_i\) is insufficient. Estimate persistence over preregistered horizons and heterogeneity across initially winning and losing positions. The broader question is whether these identified mechanisms recur with economically important magnitudes in consequential saving and investment decisions, under comparable information and ownership conditions. It should not be treated as a theorem that all optimistic beliefs are motivated or that the stock-market findings already identify population prevalence.

\par\textbf{Formulation and scope.} Proposed causal research design extending a supporting experiment. Full-paper agenda and persistence/prevalence claims could not be verified.

\par\textbf{Open-question provenance.} An explicit published open-question statement was not verified. Sources below are supporting literature and must not be treated as a first listing.

\par\textbf{Historical priority.} The source supplies evidence in stock-return experiments, not verified priority or an explicit statement of the broad saving and self-image question.

\subsection{Short English statement}
Identify how incentives for desirable conclusions change economic beliefs and actions under the specified randomized design and measurement assumptions.

\subsection{Sources}
\begin{itemize}\raggedright
\item[\textnormal{[S1]}]\hypertarget{source:30007587:1}{} Carlos Cueva, Iñigo Iturbe-Ormaetxe. 2025. Motivated beliefs about stock returns. Abstract, corroborated by author research page https://sites.google.com/site/carloscuevaherrero/home/research. Full paper was not accessible in this audit..
\url{https://doi.org/10.1016/j.jbankfin.2025.107510}.
DOI: 10.1016/j.jbankfin.2025.107510.
{\small\textit{Supporting or current literature; see evidence qualification. The authors report experimental effects of stock purchase on optimistic expectations after losses and a relationship with reluctance to sell. The source does not verify the broader saving, self-image, persistence, and prevalence question as a stated agenda.}}
\end{itemize}
\end{document}
